\documentclass{tokh}

\begin{document}

\centerline{\includegraphics[width=\linewidth]{art/cover/tokh-banner.pdf}}

\bigskip
\bigskip
\noindent{\color{TitleColor}\fontsize{36}{36}\neuropolitical\textpdfrender{TextRenderingMode=FillStroke,
    LineWidth=1.5pt,
    FillColor=White,
    StrokeColor=TitleColor,
  }{Organizers' Notes}}

\noindent%
Thank you for organizing a campaign of \textit{The Gifts of Dr. Tokh!}
This section presents a few notes to help you execute.  Importantly,
there is no secret or other decision making on the organizer's side in
this campaign, so there is no reason you can't also participate as a
player.

If you have questions about executing the campaign, please email us
via \url{contact@rocketshipgames.com}, and the resources mentioned
below are on our website at \url{http://rocketshipgames.com}.

\section{Players}

The campaign as-is should work well for any moderate size group of
players, from just two players to about 24.  Beyond that the players
probably need to be organized into sub-groups to expedite pairings.
If you think you'll be running such an event and need or want help,
email us via \url{contact@rocketshipgames.com}!  We have experience
running similar campaigns with up to~160+ participants and would be
happy to discuss techniques for managing large groups.

Groups of two or just a few players should feel free to play multiple
games per round, such that everyone gets a choice of locations and
missions.  Campaign progressions still be applied once per round
rather than per game, with experience gained and lists made
unavailable between sets of matches.

The campaign's mechanics also support players missing rounds or
joining late without much hassle.  Players dropping out won't affect
much other than potentially having an odd number of players---players
aren't uniquely tied to map locations or anything else that might
require their continued personal participation, including scoring.
New players can also be added, again just with the possibility of
having an odd number of players.  Players joining late should apply
enough campaign progressions to catch up.

In dividing into alliances, it is up to you, your players, and the
armies your group has available whether or not to align them by
factions.  The minimal backstory of both the campaign and
\textit{Infinity} in general make it eminently plausible to have
different factions teaming up or identical factions fighting.

Some effort should be spent to balance the alliances by skill as best
as possible.

As noted, the concluding Datacenter mission is designed to be a
doubles game combining RECON+ lists.  But there is little reason it
couldn't be held as a standard 300pt game if you or your players wish.
This fact may also be utilized to workaround an uneven number of
players for the final round.

\clearpage
\section{Materials}

In addition to the usual tables, terrain, and objectives, you'll need
the following materials:

\bigskip
\centerline{\begin{tabular}{|C{2in}L{5in}|}\arrayrulecolor{CornflowerBlue}
  \rowcolor{CornflowerBlue}\textcolor{White}{\textbf{Resource}} & \textcolor{White}{\textbf{Note}}\\
  \\[-6pt]
  Karma and Network Directory Tokens & Roughly~6 each per
                                       player will be needed for the campaign.\\
  \\[-6pt]
  \rowcolor{Gray!16} Bomb Markers & At most 12 are needed in each Physical Plant
                                    mission.\\
  \\[-6pt]
  Xenomorphs & Models or proxy markers, four are needed in each
               Bio/Xeno mission.\\
  \\[-6pt]
  \rowcolor{Gray!16} Datacenter Tokens & One set of 9 is needed for each pair of
                                         players.\\
  \\[-6pt]
  Scorecards & One is needed for each mission for players to report results.\\
  \\[-6pt]
  \rowcolor{Gray!16} Hobby Ballots & [optional] One per player if conducting a hobby/painting vote.\\
  \\[-6pt]
  Map \& Alliance Tokens &  [optional] To visually present campaign status.\\
\hline
\end{tabular}}

\bigskip%
PDFs to easily create all of these out of paper are available on the
website.  Some judgement will be needed to determine how many of some
of these to make, based on the number of players and the missions
they'll play, as not every match will play the same mission each
round.

Xenomorph proxy markers are included in the PDFs, but options for
easily available models include:
\begin{squishitemize}
  \item Corvus Belli: Hungries
  \item Games Workshop: Hormagaunts or Genestealers
  \item A wide variety of papercraft aliens may be found with quick Googling
\end{squishitemize}


\vspace{-9pt}
\section{Schedule}

Fairly casual RECON+ matches can be held in under~90 minutes.  Doubles
games in this format require about 2.5 hours.  The campaign is
therefore flexible to being held over a single day or several
sessions.

A reasonable schedule for a single day event is:

\bigskip
\centerline{\begin{tabular}{|C{1in}L{4in}|}\arrayrulecolor{CornflowerBlue}
  \rowcolor{CornflowerBlue}\textcolor{White}{\textbf{Time}} & \textcolor{White}{\textbf{Activity}}\\
  11:00 & Doors open\\
  \rowcolor{Gray!16}   11:55 & Registration closes and alliances determined\\
  12:00 & Campaign briefing and pairings\\
  \rowcolor{Gray!16}   12:15 & Round 1 \hfill [90 minutes]\\
  13:45 & Campaign briefing and pairings\\
  \rowcolor{Gray!16}   14:00 & Round 2 \hfill [90 minutes]\\
  15:30 & Campaign briefing and pairings\\
  \rowcolor{Gray!16}   15:45 & Round 3 \hfill [90 minutes]\\
  17:15 & Dinner break\\
  \rowcolor{Gray!16}   18:00 & Campaign briefing and pairings\\
  18:15 & Round 4 \hfill [2.5 hours]\\
  \rowcolor{Gray!16}   21:00 & Campaign briefing and awards\\
  \hline
\end{tabular}}

\bigskip%
Campaign briefings should present the current status of the map and
alliance scores, give the alliances a few minutes to discuss strategy
on their own, and then conduct pairings.

\section{Scoring}

The website has a link to a Google Docs spreadsheet workbook for
scoring the campaign and conducting pairings, with detailed
instructions.  It automates almost all of the tabulations and
calculations, requiring organizers to simply enter a player roster and
match results (and painting votes if they wish).

Otherwise, at minimum, organizers need to track:
\begin{squishitemize}
\item Total objective points earned in each map location by each alliance;
\item Total objective points earned overall by each alliance.
\end{squishitemize}

\section{Pairings}

Pairings are conducted in a very simple team-driven fashion.  At the
beginning, the Creditors have the initiative.  Subsequently the
alliance with the most total objective points earned so far has the
initiative.  In each round, beginning with the alliance holding the
initiative, they each alternate putting forward a player and a target
location.  The other alliance then responds with an opponent and a
game table.  Note that the Breach Points all use the same mission
rules, but are four separate campaign locations.  Also, everyone
fights in the \textit{Datacenter} in the final round.

This structure permits players to have a discussion within their
alliances and collectively make strategic decisions about which of
their squads and personal styles are best suited to go after
particular mission objectives or defend against attacking players.

Alliances should be highly encouraged to let each player choose a
mission in the second and third rounds at least once, rather than
defending reactively in both.

\section{Prizes}

This is intended to be a narrative campaign, but it's still possible
to award prizes.

We highly recommend a vote among the players for hobby/painting
prizes, to encourage painted models.  A simple but effective way to do
this is to have players list their top~3 favorite armies and then give
each of those armies 3, 2, and 1 vote in order.  This could be done
over armies (e.g., players' strike team cores), or individual newly
painted models to encourage continuing progress, or both.

It's reasonable to award small prizes based on game performance, but
they should be kept small and distributed well to avoid an overly
competitive environment.  We recommend equal awards for the top placed
player on each alliance, as they won't necessarily have faced other
strong players on their own alliance, and the pairings don't
necessarily match up equal pairings as they do in an tournament.

Another idea is to offer prizes for particular achievements not
strictly related to winning games, e.g., first to take down a TAG,
most Network Directory tokens collected, etc..

\end{document}
